\documentclass[11pt]{article}
\usepackage[margin=1in]{geometry}
\usepackage{xcolor}
\usepackage{hyperref}
\usepackage{booktabs}

\definecolor{revcomment}{rgb}{0.2,0.2,0.6}
\definecolor{response}{rgb}{0,0.4,0}

\newcommand{\reviewer}[1]{\noindent\textbf{\large Reviewer #1}\medskip\hrule\medskip}
\newcommand{\rcomment}[1]{\noindent\textcolor{revcomment}{\textit{Reviewer comment:}} #1\medskip}
\newcommand{\rresponse}[1]{\noindent\textcolor{response}{\textbf{Response:}} #1\medskip\medskip}

\title{Response to Reviewers\\Manuscript ID: Access-2025-56953}
\author{Gourav Singla}
\date{}

\begin{document}
\maketitle

We thank the Associate Editor (Dr. Shadi Alawneh) and both reviewers for their detailed and constructive feedback on the original submission. Their comments significantly strengthened this revised manuscript. Below we address each comment in turn. Major changes to the manuscript are summarized first, followed by point-by-point responses.

\section*{Summary of Major Revisions}

In response to reviewer feedback, the manuscript has been substantially expanded:

\begin{enumerate}
    \item \textbf{Multiple generation models.} The original paper used only GPT-4o. The revision adds \textbf{Claude Sonnet 4} and \textbf{Gemini 2.0 Flash} as generation models, totaling three independent LLMs from three different providers. Method rankings are consistent across all three (Section IV.C, Table V).
    \item \textbf{Multiple downstream tasks.} The original paper evaluated only slide generation. The revision adds a \textbf{factual question answering (Q\&A) task} with 10 questions per document and an objective correctness metric (1{,}741 evaluations, $\sim$17K judgments). Q\&A reveals dramatically larger chunking-method differences than slides (Section IV.E).
    \item \textbf{Multiple LLM judges.} The original paper used only Gemini 2.0 Flash as judge. The revision adds \textbf{GPT-4o-mini} and \textbf{Claude Sonnet 4} as additional independent judges, with cross-judge agreement analysis (Section IV.D, Table VII). Three judges from three model families produce identical method rankings.
    \item \textbf{Human evaluation.} \textbf{50 blind document pairs} were rated by an expert annotator, validating AI judge findings (Section IV.F, Table X).
    \item \textbf{Expanded dataset.} The corpus grew from 101 research papers to \textbf{218 diverse documents} including non-research material (government reports, NGO whitepapers, industry outlooks, technical manuals, educational materials) spanning 20+ domains.
    \item \textbf{Theoretical framework.} A new \textbf{Section III: Theoretical Framework} formalizes the information-density model, derives the W-curve from lost-in-the-middle theory, and predicts the conditions under which PAC should outperform truncation. These theoretical predictions are then empirically tested.
    \item \textbf{Ablation studies.} Chunking was re-run at four budgets (1K, 2K, 3K, 5K words) on all 218 documents (Section IV.H, Table XI).
    \item \textbf{W-curve sensitivity.} PAC was re-run with 11 alternative parameter configurations (Section IV.I, Table XII) showing the method is robust to parameter perturbation.
    \item \textbf{Novel finding.} Length-stratified analysis (Section IV.B, Table IV) reveals that PAC \textbf{outperforms all alternatives for medium-length documents (5K--10K words) on slide generation}---a niche missed in the original paper.
    \item \textbf{Practitioner guidelines} (Section V.C) translate findings into task-aware recommendations.
\end{enumerate}

The revised manuscript reflects approximately 6{,}500 slide evaluations, 1{,}741 Q\&A evaluations, 50 human ratings, plus the original 404 baseline evaluations---an order-of-magnitude expansion of the original empirical scope.

\bigskip
\hrule
\bigskip

\reviewer{1}

\rcomment{The contribution lacks generalizability: the study uses a single generation model (GPT-4o), a single downstream task (slide generation), and a single LLM evaluator (Gemini). It is unclear whether the findings hold beyond this narrow setup.}

\rresponse{We agree that single-model/task/judge studies risk being unrepresentative. The revised paper addresses each axis directly:
\begin{itemize}
    \item \textbf{Three generation models}: GPT-4o, Claude Sonnet 4, Gemini 2.0 Flash. Each was tested on every (document, method) pair (Section III.C). Method rankings are consistent across all three (Table V), showing the findings are not model-specific.
    \item \textbf{Two downstream tasks}: slide generation (subjective quality) and factual Q\&A (objective correctness). The two tasks reveal qualitatively different chunking sensitivities, which is itself a key contribution.
    \item \textbf{Three LLM judges}: Gemini 2.0 Flash, GPT-4o-mini, and Claude Sonnet 4---one from each major provider. All three produce identical rankings; pairwise Pearson correlations between judges are $r = 0.53$--$0.61$ (Table VII).
    \item \textbf{Human evaluation}: 50 blind pairs validate that AI judges agree with human preferences (Table X).
\end{itemize}
The cumulative evidence (3 generators $\times$ 3 judges $\times$ human, plus objective Q\&A correctness) makes the findings substantially more robust than the original submission.}

\rcomment{The paper would benefit from human evaluation. LLM-as-judge methodologies are useful but should be cross-checked.}

\rresponse{We added a human evaluation study on \textbf{50 randomly selected document pairs} (Section IV.F). For each pair, an expert annotator was shown two anonymized slide decks side-by-side---one from PAC, one from a baseline---and asked to choose which deck was better on six dimensions (completeness, accuracy, statistics retention, coherence, relevance, overall) plus a confidence score. The answer key was sealed until ratings were complete. Results (Table X) confirm AI-judge findings: PAC is preferred in only 38\% of overall comparisons. This rules out single-judge bias as the explanation for PAC's lower scores. We acknowledge in Section V.E that single-rater is a remaining limitation; multi-rater agreement is targeted for future work.}

\rcomment{The paper needs ablation studies. Without them, it is unclear which design choices drive the results.}

\rresponse{We added three ablation/sensitivity studies covering all design dimensions raised by the reviewer:
\begin{itemize}
    \item \textbf{Chunk budget ablation} (Section IV.H, Table XI): All four methods were re-run with budgets of 1{,}000, 2{,}000, 3{,}000, and 5{,}000 words on all 218 documents. Method rankings are stable across budgets, with diminishing returns above 3K words.
    \item \textbf{W-curve parameter sensitivity} (Section IV.I, Table XII): PAC was re-run with 11 alternative W-curve configurations varying intro weight, conclusion weight, results-region peak, and floor (e.g., \texttt{baseline}, \texttt{intro\_low}, \texttt{intro\_high}, \texttt{concl\_low}, \texttt{concl\_high}, \texttt{res\_low}, \texttt{res\_high}, \texttt{floor\_zero}, \texttt{floor\_high}, \texttt{flat}, \texttt{u\_curve}). Output size variance is $<$2\% across configurations, demonstrating robustness to parameter perturbation.
    \item \textbf{Prompt structure}: All generation calls used the same prompt template across methods, generators, and tasks (held constant by design to isolate the chunking variable). We acknowledge in Section V.E that prompt-aware chunking is a complementary research direction not addressed in this study.
\end{itemize}}

\rcomment{The paper lacks theoretical depth. Why should position-aware chunking work? When should it succeed or fail? For example, if truncation works well, is this because research papers are inherently front-loaded, because slide generation is high-level summarization, or because the evaluation method is insensitive to deeper content differences? Likewise, if position-aware chunking underperforms, does this suggest positional heuristics are fundamentally weak, or only that the current formulation is too simplistic?}

\rresponse{We addressed this comment in three places:
\begin{enumerate}
    \item \textbf{New Section III (Theoretical Framework)}: Formalizes chunking as a budgeted utility-maximization problem with two components: information density $I(p)$ along the document and LLM attention weight $w_{\text{LLM}}(p)$ along the context (Section III.A). The W-shaped position curve is derived from the combination of these factors (Section III.B, Equation 1). Three conditions under which PAC should outperform truncation are stated explicitly (Section III.C): long documents, rich middle content, and LLMs that integrate discontinuous chunks. These theoretical predictions are then empirically tested in Section IV.B, where we find PAC's empirical sweet spot is 5--10K-word documents---consistent with the theory.
    \item \textbf{New Section V.B (``Why Does Truncation Win? Three Hypotheses'')}: We address the reviewer's three candidate explanations directly. We use the empirical results to confirm Hypothesis 1 (front-loading), confirm Hypothesis 2 (slide generation is a high-level summarization task that does not stress fine-grained recall), and \textit{reject} Hypothesis 3 (evaluation insensitivity)---ruling it out using cross-judge agreement, human evaluation alignment, and the 18-percentage-point Q\&A gap that the same evaluation infrastructure produces.
    \item \textbf{New Section V.C (``Why Does PAC Underperform? Fundamental or Formulation?'')}: We address the reviewer's question of whether PAC's underperformance reflects fundamental weakness of positional heuristics or only simplicity of the current formulation. We argue (with evidence from the medium-document niche and the W-curve sensitivity analysis) that positional heuristics are not fundamentally weak, but that the current PAC formulation is overly simplistic in three identifiable ways: (1) document-structure assumption (no genre awareness), (2) discontinuity penalty (no coherence-aware selection), and (3) static budget allocation. Each is mapped to a future-work direction.
\end{enumerate}}

\rcomment{Single dataset (101 research papers) is too narrow. Findings may not generalize to other document types.}

\rresponse{The corpus has been expanded to \textbf{218 documents} (Section III.A, Table I). Approximately 97 are research papers (the original genre), 74 are non-research documents (government reports, NGO whitepapers, industry outlooks, corporate annual reports, technical manuals, educational lectures), and 47 are mixed (review articles, policy documents, case studies). The 20+ domains span computer science, medicine, economics, education, environmental science, engineering, law, finance, and others. Length-stratified analysis (Section IV.B) shows the findings hold across this expanded distribution.}

\rcomment{The discussion section could be improved by providing a more detailed analysis of why certain methods perform better or worse, particularly the consistent loss of statistical information during generation.}

\rresponse{Section IV.L (``Statistics Retention: A Universal Weakness'') and the discussion expand on this point. We attribute the universally poor statistics retention ($\sim$2.6/5 across all judges and all generators) to a generation-side rather than chunking-side limitation: LLMs systematically prefer to summarize trends rather than cite specific numbers, regardless of whether the input chunk contains them. This is supported by the fact that all four chunking methods show similar statistics retention scores despite very different content selection strategies---if the limitation were on the chunking side, we would expect PAC (which explicitly rewards numeric content) to retain more statistics than truncation (which has no such bias). We discuss this in Section V (``Why Does Truncation Win?'') and identify prompt-side and post-processing techniques as the relevant future-work directions in the Conclusion.}

\rcomment{The conclusions should be framed more cautiously, emphasizing that the results may be task-specific and dataset-specific, rather than general recommendations for all LLM-based content generation systems.}

\rresponse{The conclusions are now framed cautiously and explicitly tied to scope. Section V.F (``Limitations and Scope'') is expanded with eight limitations (task coverage, prompt structure, domain coverage, question generation source, single human rater, language, static budgets, fixed evaluation rubric) and ends with: ``These limitations frame the findings as task- and dataset-specific rather than universal claims.'' The practitioner guide in Section V.D explicitly conditions recommendations on document length and downstream task type rather than offering a single universal recommendation. The conclusion likewise frames the take-aways as task-aware rather than universal, and the abstract now uses phrases like ``has a clear niche'' and ``most practical applications'' to avoid overclaiming.}

\bigskip
\hrule
\bigskip

\reviewer{2}

\rcomment{Justification of W-curve parameter values is missing. Why these specific weights for intro/conclusion/results regions?}

\rresponse{The original W-curve parameters were chosen heuristically. The revised paper provides three-fold justification:
\begin{enumerate}
    \item \textbf{Theoretical derivation} (Section III.B): The W-curve is derived from the combination of information density $I(p)$ (which peaks in intro/results/conclusion for academic papers) and LLM attention $w_{\text{LLM}}(p)$ (which peaks at boundaries per Liu et al.\ 2024).
    \item \textbf{Sensitivity analysis} (Section IV.I, Table XII): We tested 11 alternative parameter configurations including \texttt{baseline}, \texttt{intro\_low/high}, \texttt{concl\_low/high}, \texttt{res\_low/high}, \texttt{floor\_zero/high}, \texttt{flat}, and \texttt{u\_curve}. Output size variance across all configurations is under 2\%, demonstrating that the qualitative shape (boundary emphasis with optional middle bump) matters more than the specific numeric values.
    \item \textbf{Empirical validation} (Section IV.B): The W-curve's predicted niche (medium-length documents with rich middle content) is empirically confirmed---PAC outperforms all baselines on slides for 5--10K-word documents.
\end{enumerate}}

\rcomment{The paper would benefit from a small human evaluation study (e.g., 20--30 pairs).}

\rresponse{We conducted human evaluation on \textbf{50 document pairs}---exceeding the suggested 20--30---in Section IV.F (Table X). Results confirm AI judge rankings.}

\rcomment{Discuss the generalizability of the position-aware chunking approach. To what kinds of documents and tasks does it apply?}

\rresponse{Section V.C (``When to Use PAC: A Practitioner Guide'') provides explicit task- and length-aware guidance:
\begin{itemize}
    \item \textbf{Short documents ($<$5K)}: Chunking method is irrelevant; budget covers most content.
    \item \textbf{Medium documents (5--10K), summarization tasks}: \textbf{Use PAC}---it outperforms all alternatives.
    \item \textbf{Long documents (10--20K), summarization tasks}: Use truncation or semantic.
    \item \textbf{Q\&A or fact extraction (any length)}: Use truncation; document beginnings concentrate factual answers.
\end{itemize}
This guidance is grounded in the length- and task-stratified results (Tables IV and IX) and replaces the original paper's monolithic ``always use X'' recommendations.}

\rcomment{A few implementation details needed for reproducibility are missing: the exact evaluation prompt given to Gemini, the chunk deduplication logic for the fixed-size approach, and the complete PAC chunk selection algorithm.}

\rresponse{Each of the three specific items is now addressed in the methodology:
\begin{enumerate}
    \item \textbf{Exact evaluation prompt}: A condensed excerpt of the multi-judge evaluation prompt now appears as a typeset block in Section IV.E (six rated dimensions, JSON output schema, document and slide-deck placeholders). The full prompt with formatting is available in the public repository.
    \item \textbf{Chunk deduplication logic for fixed-size}: We added a clarification at the end of Section IV.B noting that fixed-size first-last does not employ a deduplication step because the first 1{,}000 and last 1{,}000 words are non-overlapping by construction. We also clarify that semantic chunking does not need explicit deduplication because breakpoint detection ensures coherent, non-redundant chunks.
    \item \textbf{Complete PAC chunk selection algorithm}: Algorithm 1 in Section IV.B now provides the full pseudocode of the PAC selection procedure, including the windowing, scoring, redundancy gate (Jaccard $\tau = 0.6$), coverage-ensuring step, and position-based reordering for narrative flow.
\end{enumerate}
In addition, all code, prompts, and corpus metadata are available at the GitHub repository linked in the Data Availability statement.}

\rcomment{The 3{,}700$\times$ speed difference between semantic chunking and truncation is a strong practical argument against semantic chunking, but this implication could be stated more directly in the discussion and abstract.}

\rresponse{We made the speed argument prominent in three places:
\begin{enumerate}
    \item \textbf{Abstract}: The closing sentence of the abstract now explicitly states the measured 3{,}300$\times$ latency advantage of truncation over semantic chunking and frames it as a strong practical case for truncation as the default.
    \item \textbf{New Section IV.K (Computational Cost)}: A dedicated subsection with a latency table (Table XIII) reports per-document latency for all four methods (truncation 0.4ms, fixed-size 0.4ms, PAC 17.9ms, semantic 1{,}327.5ms), translates this into pipeline-scale impact (e.g., 3.7 hours vs 4 seconds for 10{,}000 documents), and explicitly states that the latency penalty makes semantic chunking ``a poor default choice for throughput-sensitive applications''.
    \item \textbf{New Section V.E (Practical Implications: The Speed--Quality Trade-off)}: A new discussion subsection ties the latency results to the quality results and frames PAC as a much better complexity/quality trade-off than semantic chunking (70$\times$ faster than semantic with comparable or better quality).
    \item \textbf{Conclusion}: A new paragraph summarizes the speed--quality trade-off as the second of the two practical takeaways of the paper.
\end{enumerate}}

\bigskip
\hrule
\bigskip

\section*{Closing}

We hope the revisions adequately address the reviewers' concerns. The cumulative empirical evidence---6{,}498 slide evaluations across three generators and three judges, 1{,}741 Q\&A evaluations with objective correctness, 50 human ratings, ablations on budgets and parameters, an expanded 218-document corpus spanning 20+ domains, and a theoretical framework predicting (and confirming) PAC's empirical niche---substantially strengthens the original contribution. The revised paper now reports a richer, more nuanced finding than the original: chunking method effects are \textit{task-dependent and length-dependent}, and we provide practitioners with actionable guidance for when each method should be used.

We thank the reviewers and editor for their time and constructive feedback.

\bigskip
\noindent\textit{Sincerely,}\\
Gourav Singla\\
Independent Researcher

\end{document}
